\chapter{Phương Pháp Luận \& Kiến Trúc Hệ Thống}

\section{Phương pháp luận thiết kế hệ thống}
Để xây dựng thành công hệ thống vi xử lý đa lõi 8-Core RISC-V hoạt động ổn định và tối ưu, dự án áp dụng phương pháp luận thiết kế phần cứng hệ thống theo các nguyên tắc cốt lõi sau:

\subsection{Quy trình thiết kế Top-Down kết hợp Bottom-Up}
Quy trình xây dựng hệ thống kết hợp hài hòa hai cách tiếp cận:
\begin{itemize}
    \item \textbf{Thiết kế từ trên xuống (Top-Down Design):} Xuất phát từ yêu cầu hệ thống đa nhân chia sẻ bộ nhớ và sử dụng tập lệnh RISC-V. Hệ thống được phân rã thành các khối chức năng độc lập bao gồm: Lõi đơn nhân (Core), Bộ phân xử (Arbiter), Bộ điều khiển Bus (Bus Controller) và Bộ nhớ dữ liệu dùng chung (Shared RAM). Việc xác định rõ giao tiếp (Interface) giữa các khối này là bước tiên quyết trước khi viết mã RTL.
    \item \textbf{Tích hợp từ dưới lên (Bottom-Up Integration):} Hiện thực hóa và kiểm thử kỹ lưỡng từ mức linh kiện cơ bản: Đơn vị số học logic (ALU), Tập thanh ghi (Register File), Bộ điều khiển (Control Unit), Logic cập nhật PC. Sau đó tích hợp thành Lõi xử lý đơn (Core) và trục liên kết đa nhân (Interconnect). Cuối cùng là tích hợp toàn hệ thống 8-Core và kiểm chứng bằng testbench hệ thống phức tạp.
\end{itemize}

\subsection{Mô hình hóa phần cứng ở mức RTL (RTL Modeling)}
Toàn bộ hệ thống được mô tả bằng ngôn ngữ mô tả phần cứng \textbf{Verilog HDL} (chuẩn IEEE 1364-2001). Phương pháp thiết kế RTL tập trung vào:
\begin{itemize}
    \item \textbf{Tách biệt mạch tuần tự và mạch tổ hợp:} Sử dụng sườn lên xung nhịp (\texttt{posedge clk}) để chốt dữ liệu vào các thanh ghi pipeline và bộ đếm chương trình PC. Sử dụng khối \texttt{always @(*)} và \texttt{assign} cho các mạch tổ hợp giải mã lệnh và tính toán ALU.
    \item \textbf{Thiết kế modular hóa:} Mỗi khối chức năng được đóng gói trong một file riêng biệt, giao tiếp qua các cổng dữ liệu (ports) rõ ràng, giúp dễ dàng bảo trì và kiểm thử độc lập.
\end{itemize}

\subsection{Quy trình mô phỏng và kiểm chứng dựa trên công cụ mã nguồn mở}
Kiểm chứng chức năng (Functional Verification) đóng vai trò sống còn. Phương pháp luận kiểm chứng của dự án dựa trên:
\begin{itemize}
    \item \textbf{Mô phỏng mức RTL (RTL Simulation):} Sử dụng công cụ mã nguồn mở \textbf{Icarus Verilog} để biên dịch mã nguồn và chạy các testbench. Các kịch bản kiểm thử được thiết kế để bao phủ các tình huống xung đột dữ liệu (data hazard), xung đột điều khiển (control hazard) và tranh chấp bus bộ nhớ.
    \item \textbf{Phân tích giản đồ dạng sóng (Waveform Analysis):} Xuất tệp dạng sóng định dạng VCD (Value Change Dump) và sử dụng phần mềm trực quan \textbf{GTKWave} để gỡ lỗi (debug) các hành vi phần cứng theo từng chu kỳ xung nhịp.
\end{itemize}

\subsection{Mô hình hóa chia sẻ tài nguyên và cơ chế đồng bộ}
Đối với hệ thống 8 lõi dùng chung bộ nhớ:
\begin{itemize}
    \item \textbf{Tránh deadlock và starvation:} Áp dụng thuật toán phân xử Round-Robin cho bộ Arbiter nhằm cấp quyền truy cập bus tuần hoàn, công bằng giữa các nhân.
    \item \textbf{Đồng bộ hóa phần cứng:} Hỗ trợ các chỉ thị lệnh nguyên tử (Atomic) từ nhóm mở rộng A-Extension nhằm giải quyết xung đột ghi/đọc đồng thời từ phần mềm đa nhân.
\end{itemize}

\section{Kiến trúc tập lệnh RISC-V (ISA) áp dụng trong dự án}
Trước khi đi vào thiết kế phần cứng ở mức RTL, việc hiểu rõ kiến trúc tập lệnh (Instruction Set Architecture - ISA) là bắt buộc. ISA đóng vai trò là đặc tả giao tiếp giữa phần mềm và phần cứng. Trong dự án này, hệ thống vi xử lý được xây dựng dựa trên tập chỉ thị \textbf{RV32I} kết hợp với một số chỉ thị thuộc nhóm mở rộng nguyên tử \textbf{A-Extension} (RV32IA subset).

\subsection{Tổng quan về tập chỉ thị RV32I}
RV32I là tập chỉ thị số nguyên cơ bản dạng 32-bit của RISC-V. Đây là tập chỉ thị bắt buộc và là nền tảng cho mọi bộ xử lý RISC-V. Các đặc trưng kỹ thuật chính bao gồm:
\begin{itemize}
    \item \textbf{Độ rộng địa chỉ và thanh ghi:} 32-bit cho cả thanh ghi đa năng và không gian địa chỉ bộ nhớ.
    \item \textbf{Cơ chế nạp/lưu (Load-Store Architecture):} Chỉ các chỉ thị Load (\texttt{LW}) và Store (\texttt{SW}) mới được phép truy cập bộ nhớ dữ liệu. Các phép toán toán học và logic chỉ thực hiện trực tiếp trên tập thanh ghi.
    \item \textbf{Độ dài lệnh cố định:} Mọi chỉ thị lệnh đều có độ rộng cố định là 32-bit, giúp đơn giản hóa quá trình nạp và giải mã lệnh trong pipeline phần cứng.
\end{itemize}

\subsection{Tập thanh ghi đa năng (Register File)}
Bộ xử lý tích hợp 32 thanh ghi đa năng kí hiệu từ \texttt{x0} đến \texttt{x31}. Trong đó, thanh ghi \texttt{x0} (zero) được nối cứng về mức logic 0 ở mức phần cứng và không thể thay đổi giá trị. Bảng \ref{tab:registers} dưới đây mô tả chi tiết chức năng và quy ước đặt tên (ABI) của tập thanh ghi:

\begin{table}[H]
\centering
\small
\begin{tabularx}{\textwidth}{l l X c}
\toprule
\textbf{Thanh ghi} & \textbf{Tên ABI} & \textbf{Mô tả chức năng} & \textbf{Lưu giữ} \\
\midrule
\texttt{x0} & \texttt{zero} & Luôn bằng 0 (Hằng số phần cứng) & -- \\
\texttt{x1} & \texttt{ra} & Địa chỉ trả về của hàm (Return Address) & Caller \\
\texttt{x2} & \texttt{sp} & Con trỏ ngăn xếp (Stack Pointer) & Callee \\
\texttt{x3} & \texttt{gp} & Con trỏ toàn cục (Global Pointer) & -- \\
\texttt{x4} & \texttt{tp} & Con trỏ luồng/tiểu trình (Thread Pointer) & -- \\
\texttt{x5--x7} & \texttt{t0--t2} & Thanh ghi tạm thời (Temporaries) & Caller \\
\texttt{x8} & \texttt{s0/fp} & Thanh ghi lưu giữ / Con trỏ khung (Frame Pointer) & Callee \\
\texttt{x9} & \texttt{s1} & Thanh ghi lưu giữ (Saved register) & Callee \\
\texttt{x10--x11} & \texttt{a0--a1} & Tham số truyền vào / Kết quả trả về của hàm & Caller \\
\texttt{x12--x17} & \texttt{a2--a7} & Các tham số truyền vào hàm (Arguments) & Caller \\
\texttt{x18--x27} & \texttt{s2--s11} & Các thanh ghi lưu giữ (Saved registers) & Callee \\
\texttt{x28--x31} & \texttt{t3--t6} & Các thanh ghi tạm thời (Temporaries) & Caller \\
\bottomrule
\end{tabularx}
\caption{Quy ước chức năng tập thanh ghi trong kiến trúc RISC-V}
\label{tab:registers}
\end{table}

\subsection{Các kiểu định dạng chỉ thị lệnh cơ bản}
Kiến trúc RISC-V phân chia mã lệnh 32-bit thành 6 kiểu định dạng cơ bản để tối ưu bộ giải mã:
\begin{itemize}
    \item \textbf{R-type (Register):} Dành cho các lệnh tính toán số học/logic giữa các thanh ghi nguồn (\texttt{rs1}, \texttt{rs2}) và ghi kết quả vào thanh ghi đích (\texttt{rd}).
    \item \textbf{I-type (Immediate):} Dành cho các lệnh tính toán với hằng số nạp trực tiếp 12-bit (\texttt{imm}) hoặc lệnh Load bộ nhớ.
    \item \textbf{S-type (Store):} Dành cho chỉ thị Store lưu dữ liệu từ thanh ghi nguồn vào bộ nhớ.
    \item \textbf{B-type (Branch):} Dành cho chỉ thị so sánh rẽ nhánh có điều kiện với độ lệch địa chỉ 12-bit.
    \item \textbf{U-type (Upper Immediate):} Dành cho chỉ thị nạp hằng số lớn 20-bit vào nửa cao của thanh ghi (như \texttt{LUI}, \texttt{AUIPC}).
    \item \textbf{J-type (Jump):} Dành cho chỉ thị nhảy không điều kiện (\texttt{JAL}).
\end{itemize}

\subsection{Đặc tả định dạng mã hóa và trường bit của các lệnh thực hiện}
Bộ xử lý trong dự án này hiện thực một tập con của kiến trúc tập lệnh RV32I cùng nhóm mở rộng nguyên tử A-Extension để điều khiển và đồng bộ tính toán đa lõi. Dưới đây là bảng đặc tả chi tiết mã hóa nhị phân 32-bit cho từng chỉ thị lệnh được triển khai thực tế trong mã nguồn RTL:

\subsubsection{1. Nhóm lệnh định dạng R (R-type)}
Các lệnh tính toán số học và logic thực thi trực tiếp giữa các thanh ghi nguồn (\texttt{rs1}, \texttt{rs2}) và lưu kết quả vào thanh ghi đích (\texttt{rd}). Chúng sử dụng chung mã opcode là \texttt{0110011} (7'h33). Chi tiết mã hóa các trường bit được trình bày trong Bảng \ref{tab:enc_r}:

\subsubsection{2. Nhóm lệnh định dạng I (I-type)}
Các lệnh tính toán với hằng số nạp trực tiếp (\texttt{imm}), lệnh nạp dữ liệu từ bộ nhớ (\texttt{LW}), hoặc lệnh nhảy gián tiếp (\texttt{JALR}) có độ dài hằng số là 12-bit.
\begin{itemize}
    \item Lệnh tính toán ALU sử dụng opcode \texttt{0010011} (7'h13).
    \item Lệnh đọc bộ nhớ \texttt{LW} sử dụng opcode \texttt{0000011} (7'h03).
    \item Lệnh nhảy gián tiếp \texttt{JALR} sử dụng opcode \texttt{1100111} (7'h67).
\end{itemize}
Chi tiết mã hóa được trình bày trong Bảng \ref{tab:enc_i}. Đối với hai chỉ thị dịch bit \texttt{SLLI} và \texttt{SRLI}, hằng số dịch (\texttt{shamt}) chỉ có độ rộng 5-bit (để dịch tối đa 31 bit), do đó trường \texttt{imm[11:5]} được nối cứng bằng \texttt{0000000}.

\subsubsection{3. Nhóm lệnh định dạng S (S-type)}
Định dạng dùng cho lệnh ghi dữ liệu vào bộ nhớ RAM dùng chung. Mã opcode duy nhất của nhóm là \texttt{0100011} (7'h23). Trường hằng số dịch chuyển địa chỉ (\texttt{imm}) được phân tách thành hai phần đặt ở hai vị trí khác nhau để tối ưu hóa đường dây giải mã trong CPU. Chi tiết mã hóa lệnh \texttt{SW} được mô tả trong Bảng \ref{tab:enc_s}:

\subsubsection{4. Nhóm lệnh định dạng B (Branch - B-type)}
Các lệnh rẽ nhánh có điều kiện so sánh giá trị giữa \texttt{rs1} và \texttt{rs2} để quyết định việc thay đổi địa chỉ PC. Chúng sử dụng chung mã opcode \texttt{1100011} (7'h63). Các bit của hằng số nhảy lệch được phân rã phức tạp nhằm đồng nhất với cấu trúc mạch giải mã của định dạng S. Chi tiết cấu trúc được thể hiện qua Bảng \ref{tab:enc_b}:

\subsubsection{5. Nhóm lệnh định dạng U (U-type) và J (J-type)}
Các chỉ thị làm việc với hằng số lớn nạp trực tiếp 20-bit vào nửa cao của thanh ghi (\texttt{LUI}, \texttt{AUIPC}) hoặc thực hiện lệnh nhảy không điều kiện tầm xa (\texttt{JAL}). Chi tiết cấu trúc mã hóa được thể hiện qua Bảng \ref{tab:enc_uj}:

\subsubsection{6. Nhóm lệnh nguyên tử đồng bộ (A-Extension)}
Đặc trưng thiết yếu của hệ thống 8 lõi xử lý chia sẻ RAM dữ liệu là việc đồng bộ hóa đa luồng thông qua các chỉ thị nguyên tử (Atomic). Các chỉ thị này sử dụng chung mã opcode \texttt{0101111} (7'h2F) và trường \texttt{funct3} mặc định là \texttt{010} (độ rộng dữ liệu word 32-bit). Lệnh được phân biệt bởi mã trường 5 bit \texttt{funct5 = instruction[31:27]}. Để đơn giản hóa mạch phần cứng, hai bit điều khiển tính nhất quán bộ nhớ là \texttt{aq} (bit 26) và \texttt{rl} (bit 25) được nối cứng về 0. Chi tiết mã hóa của các lệnh nguyên tử được thể hiện qua Bảng \ref{tab:enc_atomic}:

\subsection{Các bảng tổng hợp đặc tả mã hóa chỉ thị lệnh}
Dưới đây là toàn bộ các bảng đặc tả mã hóa nhị phân chi tiết cho từng nhóm chỉ thị lệnh được hiện thực hóa trong dự án:

\begin{table}[H]
\centering
\small
\begin{tabularx}{\textwidth}{l c c c c c c X}
\toprule
\textbf{Lệnh} & \makecell{\textbf{funct7} \\ \textbf{[31:25]}} & \makecell{\textbf{rs2} \\ \textbf{[24:20]}} & \makecell{\textbf{rs1} \\ \textbf{[19:15]}} & \makecell{\textbf{funct3} \\ \textbf{[14:12]}} & \makecell{\textbf{rd} \\ \textbf{[11:7]}} & \makecell{\textbf{opcode} \\ \textbf{[6:0]}} & \textbf{Mô tả hoạt động} \\
\midrule
\texttt{ADD} & \texttt{0000000} & \texttt{rs2} & \texttt{rs1} & \texttt{000} & \texttt{rd} & \texttt{0110011} & \texttt{R[rd] = R[rs1] + R[rs2]} \\
\texttt{SUB} & \texttt{0100000} & \texttt{rs2} & \texttt{rs1} & \texttt{000} & \texttt{rd} & \texttt{0110011} & \texttt{R[rd] = R[rs1] - R[rs2]} \\
\texttt{SLL} & \texttt{0000000} & \texttt{rs2} & \texttt{rs1} & \texttt{001} & \texttt{rd} & \texttt{0110011} & \texttt{R[rd] = R[rs1] << R[rs2][4:0]} \\
\texttt{SLT} & \texttt{0000000} & \texttt{rs2} & \texttt{rs1} & \texttt{010} & \texttt{rd} & \texttt{0110011} & \texttt{R[rd] = (R[rs1] < R[rs2]) ? 1 : 0} \\
\texttt{XOR} & \texttt{0000000} & \texttt{rs2} & \texttt{rs1} & \texttt{100} & \texttt{rd} & \texttt{0110011} & \texttt{R[rd] = R[rs1] \textasciicircum\ R[rs2]} \\
\texttt{SRL} & \texttt{0000000} & \texttt{rs2} & \texttt{rs1} & \texttt{101} & \texttt{rd} & \texttt{0110011} & \texttt{R[rd] = R[rs1] >> R[rs2][4:0]} \\
\texttt{OR} & \texttt{0000000} & \texttt{rs2} & \texttt{rs1} & \texttt{110} & \texttt{rd} & \texttt{0110011} & \texttt{R[rd] = R[rs1] \textbar\ R[rs2]} \\
\texttt{AND} & \texttt{0000000} & \texttt{rs2} & \texttt{rs1} & \texttt{111} & \texttt{rd} & \texttt{0110011} & \texttt{R[rd] = R[rs1] \& R[rs2]} \\
\bottomrule
\end{tabularx}
\caption{Mã hóa chi tiết các lệnh định dạng R-type}
\label{tab:enc_r}
\end{table}

\begin{table}[H]
\centering
\small
\begin{tabularx}{\textwidth}{l c c c c c X}
\toprule
\textbf{Lệnh} & \makecell{\textbf{imm[11:0]} \\ \textbf{[31:20]}} & \makecell{\textbf{rs1} \\ \textbf{[19:15]}} & \makecell{\textbf{funct3} \\ \textbf{[14:12]}} & \makecell{\textbf{rd} \\ \textbf{[11:7]}} & \makecell{\textbf{opcode} \\ \textbf{[6:0]}} & \textbf{Mô tả hoạt động} \\
\midrule
\texttt{ADDI} & \texttt{imm[11:0]} & \texttt{rs1} & \texttt{000} & \texttt{rd} & \texttt{0010011} & \texttt{R[rd] = R[rs1] + imm} \\
\texttt{SLTI} & \texttt{imm[11:0]} & \texttt{rs1} & \texttt{010} & \texttt{rd} & \texttt{0010011} & \texttt{R[rd] = (R[rs1] < imm) ? 1 : 0} \\
\texttt{XORI} & \texttt{imm[11:0]} & \texttt{rs1} & \texttt{100} & \texttt{rd} & \texttt{0010011} & \texttt{R[rd] = R[rs1] \textasciicircum\ imm} \\
\texttt{ORI} & \texttt{imm[11:0]} & \texttt{rs1} & \texttt{110} & \texttt{rd} & \texttt{0010011} & \texttt{R[rd] = R[rs1] \textbar\ imm} \\
\texttt{ANDI} & \texttt{imm[11:0]} & \texttt{rs1} & \texttt{111} & \texttt{rd} & \texttt{0010011} & \texttt{R[rd] = R[rs1] \& imm} \\
\texttt{SLLI} & \texttt{0000000} \textbar\ \texttt{shamt[4:0]} & \texttt{rs1} & \texttt{001} & \texttt{rd} & \texttt{0010011} & \texttt{R[rd] = R[rs1] << shamt} \\
\texttt{SRLI} & \texttt{0000000} \textbar\ \texttt{shamt[4:0]} & \texttt{rs1} & \texttt{101} & \texttt{rd} & \texttt{0010011} & \texttt{R[rd] = R[rs1] >> shamt} \\
\texttt{LW} & \texttt{imm[11:0]} & \texttt{rs1} & \texttt{010} & \texttt{rd} & \texttt{0000011} & \texttt{R[rd] = M[R[rs1] + imm]} \\
\texttt{JALR} & \texttt{imm[11:0]} & \texttt{rs1} & \texttt{000} & \texttt{rd} & \texttt{1100111} & \texttt{R[rd] = PC + 4; PC = (R[rs1] + imm) \& \textasciitilde 1} \\
\bottomrule
\end{tabularx}
\caption{Mã hóa chi tiết các lệnh định dạng I-type}
\label{tab:enc_i}
\end{table}

\begin{table}[H]
\centering
\small
\begin{tabularx}{\textwidth}{l c c c c c c X}
\toprule
\textbf{Lệnh} & \makecell{\textbf{imm[11:5]} \\ \textbf{[31:25]}} & \makecell{\textbf{rs2} \\ \textbf{[24:20]}} & \makecell{\textbf{rs1} \\ \textbf{[19:15]}} & \makecell{\textbf{funct3} \\ \textbf{[14:12]}} & \makecell{\textbf{imm[4:0]} \\ \textbf{[11:7]}} & \makecell{\textbf{opcode} \\ \textbf{[6:0]}} & \textbf{Mô tả hoạt động} \\
\midrule
\texttt{SW} & \texttt{imm[11:5]} & \texttt{rs2} & \texttt{rs1} & \texttt{010} & \texttt{imm[4:0]} & \texttt{0100011} & \texttt{M[R[rs1] + imm] = R[rs2]} \\
\bottomrule
\end{tabularx}
\caption{Mã hóa chi tiết lệnh định dạng S-type}
\label{tab:enc_s}
\end{table}

\begin{table}[H]
\centering
\small
\begin{tabularx}{\textwidth}{l c c c c c c X}
\toprule
\textbf{Lệnh} & \makecell{\textbf{imm[12\textbar10:5]} \\ \textbf{[31:25]}} & \makecell{\textbf{rs2} \\ \textbf{[24:20]}} & \makecell{\textbf{rs1} \\ \textbf{[19:15]}} & \makecell{\textbf{funct3} \\ \textbf{[14:12]}} & \makecell{\textbf{imm[4:1\textbar11]} \\ \textbf{[11:7]}} & \makecell{\textbf{opcode} \\ \textbf{[6:0]}} & \textbf{Điều kiện rẽ nhánh} \\
\midrule
\texttt{BEQ} & \texttt{imm[12\textbar10:5]} & \texttt{rs2} & \texttt{rs1} & \texttt{000} & \texttt{imm[4:1\textbar11]} & \texttt{1100011} & \texttt{if (R[rs1] == R[rs2]) PC = PC + imm} \\
\texttt{BNE} & \texttt{imm[12\textbar10:5]} & \texttt{rs2} & \texttt{rs1} & \texttt{001} & \texttt{imm[4:1\textbar11]} & \texttt{1100011} & \texttt{if (R[rs1] != R[rs2]) PC = PC + imm} \\
\texttt{BLT} & \texttt{imm[12\textbar10:5]} & \texttt{rs2} & \texttt{rs1} & \texttt{100} & \texttt{imm[4:1\textbar11]} & \texttt{1100011} & \texttt{if (R[rs1] < R[rs2]) PC = PC + imm} \\
\texttt{BGE} & \texttt{imm[12\textbar10:5]} & \texttt{rs2} & \texttt{rs1} & \texttt{101} & \texttt{imm[4:1\textbar11]} & \texttt{1100011} & \texttt{if (R[rs1] >= R[rs2]) PC = PC + imm} \\
\bottomrule
\end{tabularx}
\caption{Mã hóa chi tiết các lệnh định dạng B-type}
\label{tab:enc_b}
\end{table}

\begin{table}[H]
\centering
\small
\begin{tabularx}{\textwidth}{l c c c X}
\toprule
\textbf{Lệnh} & \makecell{\textbf{imm} \\ \textbf{[31:12]}} & \makecell{\textbf{rd} \\ \textbf{[11:7]}} & \makecell{\textbf{opcode} \\ \textbf{[6:0]}} & \textbf{Mô tả hoạt động} \\
\midrule
\texttt{LUI} & \texttt{imm[31:12]} & \texttt{rd} & \texttt{0110111} & \texttt{R[rd] = imm << 12} \\
\texttt{AUIPC} & \texttt{imm[31:12]} & \texttt{rd} & \texttt{0010111} & \texttt{R[rd] = PC + (imm << 12)} \\
\texttt{JAL} & \texttt{imm[20\textbar10:1\textbar11\textbar19:12]} & \texttt{rd} & \texttt{1101111} & \texttt{R[rd] = PC + 4; PC = PC + imm} \\
\bottomrule
\end{tabularx}
\caption{Mã hóa chi tiết các lệnh định dạng U-type và J-type}
\label{tab:enc_uj}
\end{table}

\begin{table}[H]
\centering
\small
\setlength{\tabcolsep}{3pt}
\begin{tabularx}{\textwidth}{l c c c c c c c c X}
\toprule
\textbf{Lệnh} & \makecell{\textbf{funct5} \\ \textbf{[31:27]}} & \makecell{\textbf{aq} \\ \textbf{[26]}} & \makecell{\textbf{rl} \\ \textbf{[25]}} & \makecell{\textbf{rs2} \\ \textbf{[24:20]}} & \makecell{\textbf{rs1} \\ \textbf{[19:15]}} & \makecell{\textbf{funct3} \\ \textbf{[14:12]}} & \makecell{\textbf{rd} \\ \textbf{[11:7]}} & \makecell{\textbf{opcode} \\ \textbf{[6:0]}} & \textbf{Mô tả hoạt động} \\
\midrule
\texttt{LR.W} & \texttt{00010} & \texttt{0} & \texttt{0} & \texttt{00000} & \texttt{rs1} & \texttt{010} & \texttt{rd} & \texttt{0101111} & \texttt{R[rd] = M[R[rs1]] (đặt khóa)} \\
\texttt{SC.W} & \texttt{00011} & \texttt{0} & \texttt{0} & \texttt{rs2} & \texttt{rs1} & \texttt{010} & \texttt{rd} & \texttt{0101111} & \texttt{R[rd] = SC(M[R[rs1]] = R[rs2])} \\
\texttt{AMOSWAP.W}& \texttt{00001} & \texttt{0} & \texttt{0} & \texttt{rs2} & \texttt{rs1} & \texttt{010} & \texttt{rd} & \texttt{0101111} & \texttt{R[rd] = M[R[rs1]]; M[R[rs1]] = R[rs2]} \\
\texttt{AMOADD.W} & \texttt{00000} & \texttt{0} & \texttt{0} & \texttt{rs2} & \texttt{rs1} & \texttt{010} & \texttt{rd} & \texttt{0101111} & \texttt{R[rd] = M[R[rs1]]; M[R[rs1]] = R[rd] + R[rs2]} \\
\texttt{AMOAND.W} & \texttt{01100} & \texttt{0} & \texttt{0} & \texttt{rs2} & \texttt{rs1} & \texttt{010} & \texttt{rd} & \texttt{0101111} & \texttt{R[rd] = M[R[rs1]]; M[R[rs1]] = R[rd] \& R[rs2]} \\
\texttt{AMOOR.W} & \texttt{01010} & \texttt{0} & \texttt{0} & \texttt{rs2} & \texttt{rs1} & \texttt{010} & \texttt{rd} & \texttt{0101111} & \texttt{R[rd] = M[R[rs1]]; M[R[rs1]] = R[rd] \textbar\ R[rs2]} \\
\texttt{AMOXOR.W} & \texttt{00100} & \texttt{0} & \texttt{0} & \texttt{rs2} & \texttt{rs1} & \texttt{010} & \texttt{rd} & \texttt{0101111} & \texttt{R[rd] = M[R[rs1]]; M[R[rs1]] = R[rd] \textasciicircum\ R[rs2]} \\
\texttt{AMOMIN.W} & \texttt{10000} & \texttt{0} & \texttt{0} & \texttt{rs2} & \texttt{rs1} & \texttt{010} & \texttt{rd} & \texttt{0101111} & \texttt{R[rd] = M[R[rs1]]; M[R[rs1]] = min(R[rd], R[rs2])} \\
\texttt{AMOMAX.W} & \texttt{10100} & \texttt{0} & \texttt{0} & \texttt{rs2} & \texttt{rs1} & \texttt{010} & \texttt{rd} & \texttt{0101111} & \texttt{R[rd] = M[R[rs1]]; M[R[rs1]] = max(R[rd], R[rs2])} \\
\texttt{AMOMINU.W}& \texttt{11000} & \texttt{0} & \texttt{0} & \texttt{rs2} & \texttt{rs1} & \texttt{010} & \texttt{rd} & \texttt{0101111} & \texttt{R[rd] = M[R[rs1]]; M[R[rs1]] = min\_u(R[rd], R[rs2])} \\
\texttt{AMOMAXU.W}& \texttt{11100} & \texttt{0} & \texttt{0} & \texttt{rs2} & \texttt{rs1} & \texttt{010} & \texttt{rd} & \texttt{0101111} & \texttt{R[rd] = M[R[rs1]]; M[R[rs1]] = max\_u(R[rd], R[rs2])} \\
\bottomrule
\end{tabularx}
\caption{Mã hóa chi tiết các lệnh nguyên tử thuộc nhóm A-Extension}
\label{tab:enc_atomic}
\end{table}

